\documentclass[10pt,a4paper]{amsart}
\usepackage[margin=19mm]{geometry}
\usepackage{amsmath,amssymb,mathtools}
\usepackage[scr=rsfs,cal=euler]{mathalfa}
\usepackage{xfrac}
\usepackage[unicode,hidelinks,bookmarks=false]{hyperref}
\newcommand{\ie}{i.e.\ }
\newcommand{\cf}{cf.\ }
\newcommand{\resp}{respectively\ }
\newcommand{\Subsection}{\subsection}
\providecommand{\nobreakdash}{\nobreak\hbox{-}}
\setlength{\emergencystretch}{2em}
\multlinegap0pt
\usepackage{amssymb}
\usepackage{enumerate}
\newtheorem{theorem}{Theorem}
\def \wt {\widetilde}
\title{The Pascal Matrix, Commuting Tridiagonal Operators and Fourier Algebras}
\author{W.~Riley Casper, Ignacio Zurri\'an}
\date{Source: arXiv:2407.21680v2 (CC BY 4.0)}
\begin{document}
\maketitle

\noindent\emph{Source and licence.} The Pascal Matrix, Commuting Tridiagonal Operators and Fourier Algebras. Version arXiv:2407.21680v2 (CC BY 4.0), distributed by arXiv under the Creative Commons Attribution 4.0 International licence, http://creativecommons.org/licenses/by/4.0/, as recorded in the arXiv licence metadata for this exact version and shown on the abstract page https://arxiv.org/abs/2407.21680v2. Full licence notice: \texttt{licenses/arxiv-2407.21680-CC-BY-4.0.txt}.\par\smallskip

\noindent\textit{Editorial scope.} Counted unit: the article's theorem T2 (source0.tex lines 727--730). The article's complete original proof of it is delivered with it (source0.tex lines 731--748, one \texttt{proof} environment of 18 lines). No mathematics is written, repaired or re-derived: the statement and the proof are the article's own lines, in the article's own notation, and no retained line is substituted or re-flowed.  No further source-local context is needed: every cross-reference of every retained line resolves inside the two delivered spans.  Nothing was omitted from any retained span, so the package carries no omission record.  No retained line contains a citation, so the wrapper carries no bibliography and no bibliography entry.  No retained line contains a figure environment, an image-inclusion command or drawing code, so no image is delivered and the package declares no asset dependency.  Editorial formatting only, in addition: an AMS \texttt{amsart} preamble that restates the article's own declarations and macros used by the retained lines, namely the environments \texttt{theorem} on separate counters, together with the standard packages those lines need.  The retained environments print in this document as Theorem 1 in order of appearance, whereas the article prints the same items with the numbering of its own full document.  The complete native source of the article as distributed in the arXiv e-print for this exact version is bundled unchanged as \texttt{sources/163158/source0.tex}.
\begin{theorem}\label{T2}
The $N\times N$ Pascal matrix $T_N$ commutes with the $N\times N$ Jacobi matrix $J_N$ whose entries are given by
$$(J_N)_{jk} = (N^2J - \wt J)_{jk},\quad 0\leq j,k < N.$$
\end{theorem}

\begin{proof}
If we define the semi-infinite matrix $J=N^2J - \wt J$, it is clear that it commutes with $T$ since $J$ and $\wt J$ commute with $T$. On the other hand, by construction, $J$ is of the form
$$J=
\left[\begin{array}{c|c}
J_N&\bf 0\\\hline
\bf 0 & *
\end{array}\right],
$$
for $J_N$ a $N\times N$ matrix. 
Also, recall that $T_N$ is just the left-upper $N\times N$ sub-matrix of $T$, i.e., 
$$T=
\left[\begin{array}{c|c}
T_N&\bf *\\\hline
\bf * & \bf*
\end{array}\right].
$$
Then, it follows that $J_N$ commutes with $T_N$.
\end{proof}

\end{document}
